\section{Data Dictionary}

\begin{longtable}{@{}p{0.25\linewidth}p{0.20\linewidth}p{0.47\linewidth}@{}}
\caption{Core research and audit fields.}\label{tab:data_dictionary}\\
\toprule
Field & Type & Definition \\
\midrule
\endfirsthead
\toprule
Field & Type & Definition \\
\midrule
\endhead
\bottomrule
\endfoot
\texttt{date} & date & US equity trading date. \\
\texttt{security\_id} & string & Persistent identifier, preferably SEC CIK based. \\
\texttt{ticker} & string & Reference-date or acquisition symbol, retained as a time-varying label. \\
\texttt{open} & float & Raw vendor opening price. \\
\texttt{high} & float & Raw vendor high price. \\
\texttt{low} & float & Raw vendor low price. \\
\texttt{close} & float & Vendor nominal closing price. Yahoo back-adjusts its historical nominal close across splits; adjusted close additionally incorporates distributions. \\
\texttt{adj\_close} & float & Vendor dividend- and split-adjusted closing price. \\
\texttt{volume} & numeric & Reported share volume in the vendor convention. \\
\texttt{dividends} & float & Cash dividend recorded on the observation date. \\
\texttt{stock\_splits} & float & Stock split factor recorded on the observation date. \\
\texttt{pct\_return} & float & Reserved research field; null in this acquisition phase. \\
\texttt{quality flags} & Boolean fields & Separate diagnostic and issue-register fields; never automatic deletion rules. \\
\texttt{source\_confidence} & string & Security-master evidence class for reference-date eligibility. \\
\end{longtable}

Raw fields preserve vendor semantics. Adjusted close is not reconstructed by
overwriting the unadjusted close. Diagnostic returns are maintained separately
from the research-facing base.
